% ============================================================================
% SECTION 1 --- INTRODUCTION
% One job per paragraph.  Every topic sentence names its relation to the one
% before it (yet / recent work does / what none of them separates / we introduce
% / because control multiplies / purpose entering last).  No paragraph is
% joined by "also".  The method paragraph carries a bold run-in and states the
% three questions and nothing else; the guarantee is its own paragraph rather
% than a subordinate clause.  No equation: eq:central lives in Sec. 3.2.
% ============================================================================
\section{Introduction}
\label{sec:intro}

Generative modeling has become a central approach to molecular design, supporting both de novo generation and the optimization of existing molecules across increasingly complex chemical spaces \citep{gomezbombarelli2018chemvae,jin2018jtvae,jin2019vjtnn,vignac2023digress,campbell2024multiflow,lee2025genmol}. Recent graph-based diffusion, flow, and discrete generative models can learn expressive molecular distributions and increasingly support conditional generation toward desired structural and functional properties \citep{vignac2023digress,campbell2024multiflow,lee2025genmol}. These advances have made molecular generation substantially more flexible, but the learned object is still typically an endpoint distribution or a generator specialized toward a particular task, rather than a reusable representation of transitions between molecular states. For molecular design, this suggests value in modeling the path between structures, since each transformation changes both the next molecule and the set of molecular states that remain reachable afterward.

Molecular editing makes this path dependence explicit, because each edit changes the transformations available from the resulting molecule. Existing editing and lead-optimization methods already exploit sequential molecular modifications through graph-to-graph generation, reinforcement learning, fragment-based search, evolutionary optimization, and MCMC \citep{jin2019vjtnn,zhou2019moldqn,jensen2019graphga,xie2021mars,fu2021mimosa}. Yet the consequence of an edit extends beyond the molecule it immediately produces. An atom insertion, bond modification, or ring operation can open or eliminate entire downstream branches of the design space. The edit with the largest immediate property gain may therefore be a poor choice under a finite design budget, while a locally unremarkable edit may create access to substantially better molecules several steps later. This distinction is consequential in lead optimization, where gains in one property can create liabilities in others, such as improving potency while compromising solubility or selectivity, and where finite evaluation budgets make it valuable to preserve alternative molecular profiles rather than repeatedly optimize the same trade-off. New measurements may also introduce requirements that were absent from the original objective, making already-reached molecules useful starting points for a revised design goal. Likewise, requirements such as preserving an important molecular feature or remaining within a predictor's trusted chemical regime may need to hold throughout optimization rather than only at the returned molecule. Molecular editing may therefore benefit from representing not only the value of individual structures, but also which molecular futures remain reachable from each state.

Despite this progress, existing approaches capture only parts of a transition-centered formulation. Learned generators can increasingly be guided toward desired properties, but their generative trajectories are often defined through corrupted, masked, or latent states rather than reusable transitions between complete molecules \citep{vignac2023digress,campbell2024multiflow,lee2025genmol}. Molecular editing and search methods instead operate directly on complete molecules, but their proposal mechanisms are typically prescribed by the optimizer and coupled to task-specific scoring \citep{jensen2019graphga,zhou2019moldqn,xie2021mars}. Recent work further shows that pretrained discrete generators can be redirected toward multiple objectives and constraints through inference-time guidance and control \citep{chen2025areuredi,pcomole2026}. Together, these advances motivate a complementary abstraction, a goal-independent generative process over molecular transitions that can be learned once and controlled as design goals change.

We introduce \textbf{\method{}} (\textbf{CO}ntrollable \textbf{M}olecular \textbf{P}rocess \textbf{O}ver \textbf{S}tochastic \textbf{E}dits), a generative framework that learns a reusable molecular transition process and controls it across design objectives (\cref{fig:ppp}). \method{} represents generation as a stochastic sequence of executable edits between complete molecules. At each molecular state, an exact executor defines what is possible by enumerating the state-dependent set of legal, trans-dimensional molecular successors. We then learn a goal-independent reference law $R_\theta$ over this same successor set from executable molecular trajectories, capturing which molecular continuations are plausible without conditioning on a downstream design goal. A separate finite-horizon controller determines what is purposeful by reweighting these molecular continuations according to the value of the futures they make reachable. Within this formulation, future-aware edit selection, exploration across competing property trade-offs, retargeting from already-reached molecules, and statewise constraints are different interventions on the same learned molecular process rather than separate task-specific generators. Because generation and control operate on the same explicit molecular support, the transition process can be learned once and redirected across design objectives without retraining the underlying generator. Our main contributions are fourfold.

% ===== GAP: HEADLINE SENTENCE =====
% One sentence with the strongest single-objective, multi-objective and
% trajectory-control results belongs here once G2/G3/G4/G7 have landed.

\begin{enumerate}[leftmargin=*,topsep=2pt,itemsep=1pt]
\item We introduce \method{}, a learned stochastic process over executable, trans-dimensional molecular rewrites between complete molecular graphs, with generation and control defined on canonical molecular successors.
\item We formulate molecular design as finite-horizon control of this learned process and derive a budget-indexed Doob $h$-transform that introduces no transitions outside the reference support and realizes exact terminal reweighting under exact future values.
\item We develop a goal-conditioned future-value controller that steers the same frozen reference process across single- and multi-objective molecular design tasks without retraining the underlying generator.
\item We demonstrate \method{} on established molecular editing and multi-objective optimization benchmarks, and show that future-aware control, multi-objective Pareto exploration, mid-trajectory retargeting, and pathwise constraints exploit capabilities of the learned molecular transition process beyond endpoint scoring.
\end{enumerate}
